\chapter{Verificación y Validación}
\label{ch:verificacion}

Este capítulo describe las actividades de \textbf{verificación y validación}
realizadas sobre el \textbf{CMOS Microscopy Payload (CMOSM)}, con el objetivo de
demostrar el cumplimiento de los requerimientos definidos en el
Capítulo~\ref{ch:requerimientos}. La verificación se ha realizado principalmente
mediante ensayos experimentales en laboratorio, análisis de datos adquiridos y
revisión del diseño e implementación del sistema.

Dado el carácter académico y de demostrador tecnológico del payload, las
actividades de verificación se enfocan en validar la funcionalidad, estabilidad y
reproducibilidad del sistema de adquisición de imágenes.

% -------------------------------------------------------------------
\section{Estrategia general de verificación}

La estrategia de verificación del CMOSM se basa en los siguientes principios:

\begin{itemize}
    \item Verificación incremental de subsistemas (sensor, iluminación,
    control).
    \item Ensayos experimentales reproducibles en entorno de laboratorio.
    \item Uso de métricas cuantitativas simples para evaluar calidad de imagen.
    \item Documentación completa de parámetros de captura y resultados.
\end{itemize}

Los métodos de verificación empleados incluyen:
\textbf{Test (T)}, \textbf{Analysis (A)}, \textbf{Inspection (I)} y
\textbf{Review (R)}, conforme a lo definido en el
Capítulo~\ref{ch:requerimientos}.

% -------------------------------------------------------------------
\section{Verificación del sistema de adquisición}

\subsection{Inicialización y control del sensor CMOS}

Se verificó el correcto funcionamiento del sensor CMOS OV5647 mediante:
\begin{itemize}
    \item Inicialización consistente del sensor en múltiples ciclos de encendido.
    \item Configuración manual de exposición y ganancia.
    \item Desactivación de mecanismos automáticos (AE y AWB).
\end{itemize}

Los ensayos confirmaron la estabilidad de los parámetros de adquisición entre
capturas consecutivas, cumpliendo los requerimientos funcionales asociados.

% -------------------------------------------------------------------
\subsection{Verificación de resolución y formato}

El sistema fue configurado para capturar imágenes a resolución nativa de
2592×1944 píxeles. Se verificó visualmente y mediante inspección de metadatos que
la resolución se mantiene constante en todos los modos de captura.

Este ensayo valida el requerimiento FR-01.

% -------------------------------------------------------------------
\section{Calibración básica del sistema}

\subsection{Dark frame}

Se adquirieron imágenes sin iluminación para caracterizar el ruido de fondo del
sensor. Estas capturas permitieron:
\begin{itemize}
    \item Identificar componentes de ruido térmico.
    \item Estimar el nivel base de señal del sensor.
\end{itemize}

Los dark frames se utilizan posteriormente para la corrección de las imágenes
adquiridas.

\subsection{Flat field}

Se realizaron capturas bajo iluminación uniforme con el objetivo de caracterizar
la no uniformidad espacial del sistema de iluminación y del sensor. Estas
imágenes permiten corregir variaciones de intensidad entre píxeles.

\subsection{Corrección aplicada}

La corrección básica aplicada a las imágenes sigue la relación:

\[
I_{corr} = \frac{I - D}{F}
\]

donde $I$ es la imagen adquirida, $D$ el dark frame y $F$ el flat field
normalizado.

% -------------------------------------------------------------------
\section{Verificación del sistema de iluminación}

El sistema de iluminación basado en pantalla OLED fue evaluado en términos de:
\begin{itemize}
    \item Estabilidad temporal de intensidad.
    \item Repetibilidad de patrones mostrados.
    \item Capacidad de conmutación rápida entre patrones.
\end{itemize}

Los ensayos mostraron una variación de intensidad inferior al 5\% entre capturas
consecutivas bajo condiciones nominales, cumpliendo los requerimientos de
estabilidad luminosa.

% -------------------------------------------------------------------
\section{Ensayos con muestras biológicas}

\subsection{Muestra de epidermis de cebolla}

Como caso de estudio representativo, se utilizaron muestras de epidermis de
cebolla, seleccionadas por su estructura celular plana y dimensiones conocidas.
Las muestras fueron colocadas en contacto cercano con el plano del sensor CMOS.

Las imágenes adquiridas permiten identificar claramente paredes celulares y
regiones internas, validando la capacidad del sistema para capturar estructuras
biológicas relevantes.

\subsection{Influencia del acoplamiento óptico}

Se evaluó cualitativamente la influencia del acoplamiento óptico entre la muestra
y el sensor. Los ensayos demostraron que la reducción de la distancia
muestra--sensor incrementa significativamente el contraste y la definición de
bordes, confirmando la importancia de este parámetro en microscopía de contacto.

% -------------------------------------------------------------------
\section{Evaluación de técnicas computacionales}

\subsection{Fourier Ptychographic Microscopy (FPM)}

Se realizaron ensayos experimentales utilizando secuencias de iluminación
estructurada para evaluar la viabilidad de reconstrucción FPM. Los resultados
obtenidos indican que, bajo las condiciones actuales de acoplamiento óptico y
estabilidad mecánica, la mejora obtenida es limitada.

Estos resultados confirman el carácter experimental de este modo y justifican su
no inclusión como mecanismo primario de formación de imagen.

\subsection{Super--resolución basada en aprendizaje profundo}

Se aplicaron técnicas de super--resolución basadas en redes neuronales sobre
imágenes adquiridas en modo nominal. El procesamiento permitió mejorar la
legibilidad visual de bordes y estructuras ya presentes en la imagen original,
sin introducir información física adicional.

% -------------------------------------------------------------------
\section{Resumen de cumplimiento de requerimientos}

La Tabla~\ref{tab:verificacion-resumen} presenta un resumen cualitativo del estado
de verificación de los principales requerimientos del payload.

\begin{table}[H]
\centering
\begin{tabular}{p{3cm} p{7cm} p{2.5cm}}
\toprule
\textbf{Requerimiento} & \textbf{Método} & \textbf{Estado} \\
\midrule
FR-01 & Test / Review & Cumplido \\
FR-02 & Test / Inspection & Cumplido \\
PR-01 & Analysis / Test & Parcial \\
ER-01 & Analysis / Test & Cumplido \\
OPR-01 & Test & Cumplido \\
\bottomrule
\end{tabular}
\caption{Resumen del estado de verificación del CMOSM.}
\label{tab:verificacion-resumen}
\end{table}

% -------------------------------------------------------------------
\section{Conclusión de la verificación}

Las actividades de verificación realizadas demuestran que el CMOSM cumple con los
requerimientos funcionales y operativos definidos para su fase actual de
desarrollo. El sistema ha sido validado como demostrador tecnológico de
microscopía sin lentes en entorno de laboratorio, identificándose claramente los
límites y oportunidades de mejora para futuras iteraciones.
